\begin{corollary}[Inversion ruler for the linearized effective defect \texorpdfstring{$\varepsilon_{\mathrm{eff}}$}{eps_eff}]\label{cor:arity-335-eps-eff-inversion}
Following the spectral gap notation of Remark \ref{rem:arity-335-eps-vs-gap},
\(
\delta_{\mathrm{spec}}=1-\rho_{3,3,5}/\lambda
\)
and the root-of-unity twist notation \(\omega_p=e^{2\pi i/p}\).
In a simplest near-coboundary linearization model, suppose the ``bad-edge residual'' after twisting the third axis (collision cocycle) is primarily contributed by the phase \(\omega_p^{\pm1}\), and its effective weight under a natural edge measure is \(\varepsilon_{\mathrm{eff}}\in[0,1]\).
Then the corresponding leading-mode gap satisfies the first-order approximation
$$
\boxed{\
1-\frac{\rho_{3,3,p}}{\lambda}
\ \approx\
\varepsilon_{\mathrm{eff}}\bigl(1-\cos(2\pi/p)\bigr)\ }.
$$
Therefore, the observed spectral gap allows one to invert an \emph{effective defect ruler}
$$
\boxed{\
\varepsilon_{\mathrm{eff}}(p)
\ :=\
\frac{1-\rho_{3,3,p}/\lambda}{1-\cos(2\pi/p)}\ }.
$$
For the \(((3,3,5))\) output with \(p=5\) (see the secondary summary), one obtains
$$
\varepsilon_{\mathrm{eff}}(5)\approx \frac{0.049378911127}{1-\cos(2\pi/5)}
\approx 0.071461831514.
$$
This numerical value lies between the uniform edge-proportion defect \(\varepsilon\lesssim 0.066667\) and the equilibrium-weighted defect \(\varepsilon_\pi\approx 0.076393\) (secondary summary), indicating the mechanism direction that ``the spectral gap is closer to the weighted defect than to the uniform counting defect'': the bad edges may be more concentrated on channels with higher Perron equilibrium measure weight.
\end{corollary}

\begin{remark}[\texorpdfstring{$\varepsilon$}{eps}-law comparison prediction: forming a dual certificate with the \texorpdfstring{$\kappa$}{kappa}-law]\label{rem:arity-335-eps-law-prediction}
If one takes the best-found uniform defect \(\varepsilon\) as a coarse proxy for \(\varepsilon_{\mathrm{eff}}\), then the above linearization model yields a comparison prediction law that does not use the \(p=5\) spectral fitting constant:
$$
\boxed{\
\frac{\rho_{3,3,p}}{\lambda}
\ \approx\
1-\varepsilon\bigl(1-\cos(2\pi/p)\bigr)\ }.
$$
Substituting \(\varepsilon\approx 1/15\) from the secondary summary gives
$$
\rho_{3,3,7}/\lambda\approx 0.9748993201,\qquad
\rho_{3,3,11}/\lambda\approx 0.9894169022,\qquad
\rho_{3,3,13}/\lambda\approx 0.9923637350.
$$
Compared with the quadratic-law prediction from inverting \(\kappa\) via \(p=5\) in the secondary summary
\(
(\rho_{3,3,7}/\lambda,\rho_{3,3,11}/\lambda,\rho_{3,3,13}/\lambda)\approx
(0.974807,0.989798,0.992695)
\),
both are at the same order of magnitude but give different systematic drift directions.
Since the spectral scan of \(((3,3,p))\) in this paper has output reproducible true values (see Table \ref{tab:real_input_40_arity_335_n2_selection_law_primes}), one can compare the fitting quality of the two certificates using the errors
\(
\Delta_p^{(\varepsilon)}:=\rho_{3,3,p}^{(\varepsilon)}/\lambda-\rho_{3,3,p}/\lambda
\)
and
\(
\Delta_p^{(\kappa)}:=\rho_{3,3,p}^{(\kappa)}/\lambda-\rho_{3,3,p}/\lambda
\):
for \(p=11,13\), the \(\varepsilon\)-law error is smaller; while for \(p=7\) the two are comparable (numerical difference at the \(10^{-3}\) level).
This provides a falsifiable discrimination interface for ``uniform counting defect vs.\ equilibrium-weighted defect'': as \(p\) increases, if the spectral data more closely follows the \(\varepsilon\)-law, then the uniform defect already captures the main mechanism; if it more closely follows the \(\kappa\)-law or its variance limit \(\kappa_\infty\), then the spectrally effective defect is closer to the weighted defect (or requires higher-order angular corrections).
\end{remark}

\begin{remark}[Third certificate: SVD low-rank ``non-rank-one energy'' \texorpdfstring{$\eta_{\mathrm{svd}}$}{eta_svd}]\label{rem:arity-335-eta-svd}
Following the flattened matrix \(T\in\RR^{5\times 9}\) of Proposition \ref{prop:arity-335-low-rank} and its singular values \(s_1^{\mathrm{raw}}\ge\cdots\ge s_5^{\mathrm{raw}}\) (energy proportions given in the secondary summary).
Define a purely scalar ``non-rank-one energy'' certificate
$$
\boxed{\
\eta_{\mathrm{svd}}
\ :=\
1-\frac{(s_1^{\mathrm{raw}})^2}{\sum_{i=1}^5 (s_i^{\mathrm{raw}})^2}\ }.
$$
From the secondary summary output
\(
(s_1^{\mathrm{raw}})^2/\sum (s_i^{\mathrm{raw}})^2\approx 0.935619
\),
one obtains
$$
\eta_{\mathrm{svd}}\approx 0.064381.
$$
This quantity is of the same order as the combinatorial defect \(\varepsilon\approx 0.066667\) and the spectrally inverted \(\varepsilon_{\mathrm{eff}}(5)\approx 0.071462\), pointing to a mechanistic interpretation:
when the third-axis cocycle is close to a coboundary, the constant tensor more closely approximates a rank-one outer product shape along the \(c\) axis, so \(\eta_{\mathrm{svd}}\) can serve as a ``constant-level drift certificate'' for the near-coboundary property, complementing \(\varepsilon\) (combinatorial certificate) and \(1-\rho/\lambda\) (spectral certificate) and amenable to parallel tracking.
\end{remark}
